\chapter{Regulatory Context}
\label{ch:regulatory}
\chaptermark{Regulatory context}

Every Strat's code runs inside a regulatory frame.  A pricer that
returns a fair value also returns numbers that feed capital
ratios, stress-test submissions, transaction reports, and an
auditor's workpaper.  A Strat who does not know which regulation
is generating the deadline on the calendar, or why a risk-engine
input is flagged ``non-modellable,'' spends the first year of the
job nodding at acronyms in meetings.

This chapter is a vocabulary tour.  The aim is to make four
acronyms (Basel~III, FRTB, CCAR, MiFID~II) intelligible the
first time a senior Strat uses them in a sentence, and to give a
crisp account of the LIBOR-to-SOFR transition, the largest
rates-market project of the 2020s, which still shapes every curve
in fixed income (\cref{ch:fixed_income_foundations}).  We aim at
literacy, not at computing regulatory capital.

\begin{backgroundbox}[Prerequisites]
From earlier chapters:
\begin{itemize}
  \item \Cref{ch:kelly_risk}: VaR and Expected Shortfall.  FRTB's
    move from VaR to ES presupposes the reader knows ES.
  \item \Cref{ch:fixed_income_foundations}: zero curves, the USD
    SOFR curve, and unsecured vs secured overnight rates.  The
    LIBOR-to-SOFR transition sits on top of that foundation.
  \item \Cref{ch:risk_production}: xVA and capital charges.  KVA
    (capital valuation adjustment) is meaningful only once the
    reader knows what ``capital'' regulators mean.
\end{itemize}
This chapter introduces no new mathematics.  It is a glossary:
acronyms expanded, deadlines dated, operational meaning explained.
\end{backgroundbox}


%% ================================================================
\section{Basel III: capital adequacy}
\label{sec:basel3}
%% ================================================================

The \textbf{Basel Accords} are international banking regulations
written by the Basel Committee on Banking Supervision at the Bank
for International Settlements (BIS) in Basel, Switzerland.  They
are not themselves law: they are \emph{standards} which
member-jurisdiction regulators (US Fed, European Banking
Authority, UK PRA, Swiss FINMA) transpose into national
rulebooks.  Banks comply with the local transposition.

\textbf{Basel~III} is the third iteration, published in 2010 as
the response to the 2007--2008 global financial crisis.  Basel~I
(1988) introduced credit-risk capital ratios.  Basel~II (2004)
added a market-risk framework and is widely regarded as having
failed the crisis: capital buffers were too thin, too low-quality,
too weakly enforced.  Basel~III raised the required capital,
tightened what counts as capital, and added two new liquidity ratios.

\subsection{The three pillars}

Basel frameworks are organised around \emph{three pillars}.  The
language is standard and appears in every Basel document.

\begin{enumerate}
  \item \textbf{Pillar~1: minimum capital requirements.}  A
    numerical rulebook: compute risk-weighted assets (RWA) from
    credit, market, and operational exposures; hold capital of
    at least a specified fraction of RWA.
  \item \textbf{Pillar~2: supervisory review.}  The local
    regulator examines the bank's internal capital assessment
    (the ICAAP in Europe) and may impose bank-specific add-ons
    over the Pillar~1 minimum.
  \item \textbf{Pillar~3: market discipline.}  The bank
    publishes standardised disclosures (the ``Pillar~3 report'')
    so that counterparties, investors, and analysts can price
    the bank's risk.  Transparency lets the market discipline
    banks that take on too much.
\end{enumerate}

\begin{keyfactbox}[Basel III three pillars]
\begin{itemize}
  \item \textbf{Pillar 1}: quantitative minimum capital ratios.
    Credit-, market-, and operational-risk RWA summed, with
    capital as a percentage of the total.
  \item \textbf{Pillar 2}: the regulator's qualitative overlay.
    Internal capital assessment, stress testing, bank-specific
    add-ons.
  \item \textbf{Pillar 3}: public disclosure.  Standardised
    tables on capital composition, RWA, and key risk metrics.
\end{itemize}
Every regulatory conversation about capital fits into one of
these pillars; the pillar tells the reader whether a requirement
is a hard ratio, a bilateral judgement, or a disclosure obligation.
\end{keyfactbox}

\subsection{Key ratios}

The point of holding regulatory capital is loss absorption: when
the bank's trades, loans, or operations lose money, equity and
equity-like instruments take the hit before depositors or
counterparties are touched.  A higher capital ratio means more
losses can be absorbed before the bank fails.  Basel~III
tightened three families of ratio; the acronyms below appear on
every bank's quarterly investor presentation.

\begin{itemize}
  \item \textbf{Common Equity Tier~1 (CET1) ratio.}  CET1
    capital is the highest-quality regulatory capital: common
    shares, retained earnings, and little else.  The CET1
    ratio is CET1 divided by RWA.  Basel~III requires a
    minimum of $4.5\%$ plus a capital conservation buffer of
    $2.5\%$, giving an effective floor near $7\%$ before any
    additional buffers.
  \item \textbf{Tier~1 capital ratio.}  Tier~1 capital is CET1
    plus Additional Tier~1 (certain loss-absorbing hybrid
    instruments).  The minimum Tier~1 ratio is $6\%$ of RWA.
  \item \textbf{Total capital ratio.}  Tier~1 plus Tier~2
    (subordinated debt and similar), with a minimum of $8\%$
    of RWA.  Tier~2 absorbs losses only in liquidation.
  \item \textbf{Leverage ratio.}  Tier~1 capital divided by
    total exposure (\emph{not} risk-weighted).  The minimum is
    $3\%$.  The leverage ratio exists because RWA-based ratios
    can be gamed by risk-weighting optimisation; a non-risk
    ratio is a backstop.
  \item \textbf{Liquidity Coverage Ratio (LCR).}  High-quality
    liquid assets divided by net cash outflow over a 30-day
    stress scenario.  Minimum $100\%$.  LCR says: ``survive a
    month of stress without selling illiquid assets.''
  \item \textbf{Net Stable Funding Ratio (NSFR).}  Available
    stable funding divided by required stable funding over a
    one-year horizon.  Minimum $100\%$.  NSFR says: ``do not
    fund long-dated assets with overnight money.''
\end{itemize}

The LCR and NSFR are the Basel~III innovations most directly
caused by 2008, when several firms had plenty of notional capital
but ran out of cash in days.  Capital and liquidity are different
problems; the post-crisis framework treats them separately.

\begin{intuitionbox}
Basel~III's slogan: ``more capital, better capital, more
liquidity.''  More: higher minimum ratios.  Better: a larger
fraction must be common equity, not hybrid instruments that
failed to absorb losses when needed.  More liquidity: LCR and
NSFR pin down cash survival at 30 days and funding stability at
one year.
\end{intuitionbox}


%% ================================================================
\section{FRTB: the Fundamental Review of the Trading Book}
\label{sec:frtb}
%% ================================================================

\textbf{FRTB} stands for the \emph{Fundamental Review of the
Trading Book}.  It is the BCBS market-risk capital standard that
replaces the Basel~II.5 framework (itself a 2009 patch on
Basel~II).  FRTB was finalised in January 2019; the original 2022
go-live slipped with pandemic delays, and jurisdictional
implementations are staggered 2024--2026.

Where Basel~III is cross-cutting (credit, market, and operational
risk), FRTB is narrowly about \emph{market risk}: capital a bank
holds against its trading-book positions moving adversely.
Market-risk capital is the number that most directly bites a
derivatives Strat, because it is computed from their pricers'
sensitivities.  FRTB did two things: it replaced
Value-at-Risk with Expected Shortfall as the risk measure, and
it drew a much harder line between the trading book (marked to
market daily) and the banking book (held to maturity) to close
long-standing arbitrage between the two.

\subsection{Trading book versus banking book}

Banks carve the balance sheet in two.  The \textbf{banking book}
holds assets the bank intends to hold to maturity, most
obviously loans.  The \textbf{trading book} holds positions to be
traded: market-making inventory, hedges, proprietary positions.
Banking-book assets attract credit-risk capital; trading-book
positions attract market-risk capital.  A long-standing problem
was regulatory \emph{arbitrage} between the two: a bank could
park a position in whichever book gave a lower charge.  FRTB
imposes hard rules on book classification and restricts movement
between them.

\subsection{IMA and SA}

FRTB offers two capital approaches per desk:

\begin{itemize}
  \item \textbf{Internal Models Approach (IMA).}  The bank uses
    its own risk model, subject to regulatory approval.
    Approval is granted \emph{per desk}: each desk passes a
    battery of model-performance tests, and approval can be
    withdrawn if the desk fails backtesting.
  \item \textbf{Standardised Approach (SA).}  A prescribed
    sensitivities-based formula: deltas, vegas, and curvatures
    per risk factor, with regulator-set weights and
    correlations.  No model needed, but capital is generally
    higher than IMA on a well-run desk.
\end{itemize}

Every bank computes SA for every desk.  IMA is a premium option:
lower capital in exchange for the cost of building and defending
internal models to the regulator.

\subsection{From VaR to Expected Shortfall}

The most visible change from Basel~II.5 to FRTB is the
\emph{risk measure}.  The old framework used \textbf{Value-at-Risk}
at $99\%$ over a 10-day horizon.  FRTB replaces this with
\textbf{Expected Shortfall} (\cref{ch:kelly_risk}) at $97.5\%$,
over a liquidity-horizon-weighted 10-day basis.

The change is conceptual.  VaR at $99\%$ is a quantile: ``the
smallest loss I exceed with probability at most $1\%$,'' silent
on tail severity.  ES at $97.5\%$ is the \emph{expected} loss
given that the loss exceeds the $97.5\%$ quantile, a tail
mean.  For a Gaussian loss distribution the two numbers are
similar; for fat tails, the case that matters, ES is
substantially larger, and the capital charge scales with tail
severity.  FRTB picked $97.5\%$ rather than $99\%$ so that on
Gaussian distributions the new and old numbers are comparable,
while on fat tails the new number is strictly more conservative.

\subsection{Liquidity horizons}

A second FRTB innovation is the \emph{liquidity horizon} overlay.
Not all risk factors can be hedged in 10 days: a one-week FX
position can, an illiquid exotic credit correlation trade cannot.
FRTB assigns each risk factor a liquidity horizon from a menu
($10$, $20$, $40$, $60$, $120$ days) and scales the shock
accordingly, so illiquid risks attract capital as if held
unhedged for their full horizon.

\subsection{Non-modellable risk factors}

The most operationally significant piece of FRTB, and the one
Strats hear about first, is the \textbf{non-modellable risk
factor (NMRF)} regime.  A risk factor counts as \emph{modellable}
under IMA only if it has real observable price history: broadly,
24 or more observable prices in the preceding year, with no gap
longer than a month.  Anything failing that test is
\emph{non-modellable}.

NMRFs do \emph{not} enter the ES calculation.  Instead they
attract a separate punitive stressed-capital charge, computed per
NMRF with no diversification benefit across NMRFs.  An IMA desk
running many NMRFs has its capital dominated by the NMRF charge
and derives little benefit from the internal model.  Keeping risk
factors modellable becomes an ongoing infrastructure discipline:
Strats negotiate with market-data teams over whether an extra
vendor feed tips a factor into modellable status and whether the
cost is worth the capital saving.

\begin{keyfactbox}[FRTB key changes]
\begin{itemize}
  \item \textbf{ES replaces VaR.}  Expected Shortfall at
    $97.5\%$ confidence, not VaR at $99\%$.  Captures tail
    severity, not just a quantile.
  \item \textbf{Per-desk IMA approval.}  Internal-model
    approval is granted and revoked desk by desk, not
    bank-wide.  A failed backtest can push a desk onto the
    standardised approach overnight.
  \item \textbf{Liquidity horizons.}  Shocks scale with how
    long it takes to hedge a given risk factor; illiquid
    factors see 40--120-day shocks rather than 10.
  \item \textbf{NMRF.}  Risk factors with insufficient price
    observations attract punitive stressed-capital charges
    outside the internal model.  Modellability becomes an
    infrastructure concern.
  \item \textbf{Trading-banking book boundary.}  Hard rules
    on book classification and movement between books, to
    close the old arbitrage.
\end{itemize}
\end{keyfactbox}

\begin{intuitionbox}
FRTB makes illiquid and data-poor risks expensive.  The
consequence is a bias towards \emph{liquid} risk factors: prefer
the vanilla swap curve over a bespoke basis spread, listed-option
vols over OTC exotic correlations, deep-market tenors over sparse
ones.  Either a desk concentrates risk in liquid buckets (lower
NMRF burden, lower capital) or accepts a structurally higher
capital footprint for the exotics franchise.  The regime is a
deliberate financial incentive to centralise risk in liquid
instruments.
\end{intuitionbox}


%% ================================================================
\section{CCAR: the Fed stress test}
\label{sec:ccar}
%% ================================================================

\textbf{CCAR} is the \emph{Comprehensive Capital Analysis and
Review}, the Federal Reserve's annual stress test for large US
bank holding companies (BHCs), introduced in 2011.  Basel is
international; CCAR is US-specific.  Both ask whether a bank has
enough capital to survive a bad scenario, but differ in mechanism.

CCAR applies to BHCs above a total-consolidated-assets threshold;
since the 2018 EGRRCPA reforms this sits around \$100 billion.
The largest US banks and the US subsidiaries of large foreign
banks are in scope; Goldman Sachs is a CCAR firm.

\subsection{The scenario-based mechanism}

Early each year the Fed publishes three scenarios: baseline,
adverse, severely adverse.  Each is a path for several dozen
macro and market variables (GDP, unemployment, the Dow, the
10-year Treasury, VIX, house prices, credit spreads) over a
nine-quarter horizon.  The severely adverse scenario is
deliberately deeper than the historical norm: a plausible
worst-case recession.

Each in-scope bank is required to:
\begin{enumerate}
  \item Project its balance sheet, revenue, losses, and
    capital ratios quarter by quarter under each Fed-mandated
    scenario.
  \item Submit those projections along with a capital plan
    (dividends, buybacks, new issuance).
  \item Show that, under the severely adverse scenario, the
    firm's post-stress capital ratios remain above the
    regulatory minima.
\end{enumerate}

The Fed also runs its own \textbf{supervisory stress test}
(DFAST, the Dodd-Frank Act Stress Test) using Fed-built models.
CCAR is the bank's exercise; DFAST is the Fed's.  Together they
form the annual stress-test cycle.

\subsection{The Strat's role}

For a markets-division Strat, the visible piece is the
trading-book loss projection.  The market-risk engine is re-run
under each Fed scenario to produce a \textbf{stressed P\&L}: what
the trading book loses, path by path, quarter by quarter.  This
feeds the bank's projected losses and through them the projected
capital ratios.

Contributing to CCAR is a recurring annual project on most
trading-desk Strat teams: scenarios land, the team maps Fed macro
variables onto pricing inputs, pricers re-run, results are
aggregated, reconciled, and submitted.  If the Fed disagrees, the
bank either reconciles or accepts the Fed's DFAST numbers.

Failing CCAR is expensive.  Before 2020 the Fed could issue
qualitative objections blocking a bank's planned capital
distributions for a year.  Since 2020 the framework has evolved
into the \textbf{Stress Capital Buffer}, in which CCAR stress
losses feed directly into the required capital buffer, making the
test an ongoing constraint rather than an annual pass/fail.

\begin{keyfactbox}[CCAR at a glance]
\begin{itemize}
  \item US-specific Federal Reserve stress test, annual, in
    scope for BHCs above roughly \$100 billion in total assets.
  \item Fed publishes baseline, adverse, and severely adverse
    macro scenarios.  Banks project nine quarters of
    balance-sheet, P\&L, and capital under each.
  \item Trading-book losses are generated by re-running the
    market-risk engine with desk-level pricer output; Strats
    own the inputs that feed the submission.
  \item Post-2020: stress losses feed the Stress Capital
    Buffer directly, making the test a continuous capital
    constraint rather than an annual pass/fail.
\end{itemize}
\end{keyfactbox}


%% ================================================================
\section{MiFID II: European market structure}
\label{sec:mifid2}
%% ================================================================

\textbf{MiFID~II} is the Markets in Financial Instruments
Directive~II, the EU's wholesale revision of securities-market
regulation.  It came into force on 3 January 2018, replacing the
original MiFID (2007).  Unlike Basel and FRTB (capital), MiFID~II
is about \emph{market structure}: how venues operate, how trades
execute, what must be reported.

Four pieces show up most often in a Strat's work.

\begin{itemize}
  \item \textbf{Pre- and post-trade transparency.}  Venues and
    systematic internalisers publish pre-trade quotes and
    post-trade execution details, with narrow
    instrument-specific deferral windows for large or illiquid
    trades.  The goal: make European markets as transparent as
    US equity markets.
  \item \textbf{Best execution.}  Firms must take \emph{all
    sufficient steps} to obtain the best result for clients
    across price, costs, speed, likelihood of execution, and
    settlement.  Annual execution-quality reports (RTS 27 and
    RTS 28) are required, subsequently scaled back by the
    ``Quick Fix'' amendments.
  \item \textbf{Research unbundling.}  Historically sell-side
    research was paid implicitly out of trading commissions.
    MiFID~II requires asset managers to pay explicitly, either
    from P\&L or a Research Payment Account funded by
    disclosed client charges.  The effect was a sharp
    contraction in sell-side research and a shift toward the
    largest names.
  \item \textbf{Transaction reporting.}  Every in-scope
    transaction is reported T+1 with typically 65 fields per
    trade, including LEIs, buyer/seller designations, and
    decision-maker identifiers.  The reporting systems are a
    significant engineering surface for any EU-facing sell-side
    firm.
\end{itemize}

MiFID~II also introduced the \textbf{Systematic Internaliser
(SI)} regime in modern form: a firm dealing on its own account,
on an organised, frequent, and systematic basis, outside a
regulated venue.  SIs have quoting and transparency obligations;
the SI designation is the main legal status under which large
dealers internalise client flow in Europe.

MiFID~II is directive-plus-regulation: the MiFIR regulation is
directly applicable EU law, while the MiFID~II directive is
transposed by each member state.  Post-Brexit the UK retains a
largely parallel regime with its own tweaks.

\begin{keyfactbox}[MiFID II essentials]
\begin{itemize}
  \item EU regulation, 3 January 2018 go-live.  Scope is
    market structure and conduct, not capital.
  \item Transparency, best execution, research unbundling,
    and T+1 transaction reporting are the four pillars most
    visible to a markets-division Strat.
  \item Introduces the Systematic Internaliser regime as the
    primary legal status for bilateral dealer flow.
\end{itemize}
\end{keyfactbox}


%% ================================================================
\section{LIBOR to SOFR}
\label{sec:libor_sofr}
%% ================================================================

For decades the most important interest-rate benchmark in global
finance was \textbf{LIBOR} (the London Interbank Offered Rate),
a daily-published set of rates intended to represent the cost
at which major banks could borrow unsecured from one another in
London, across five currencies (USD, GBP, EUR, JPY, CHF) and
seven tenors (overnight out to 12 months).  At its peak LIBOR was
embedded in an estimated \$300 trillion of contracts: swaps,
floating-rate notes, mortgages, student loans, corporate loans,
structured products.

\subsection{Why LIBOR was replaced}

LIBOR was produced by a daily \emph{submission} from a small
panel of banks.  Each submitted the rate it believed it could
borrow at; the top and bottom quartiles were trimmed and the mean
of the rest published as that day's LIBOR per currency and tenor.
Submissions were not tied to observable transactions, because
true unsecured interbank borrowing had become thin after 2008.
LIBOR was largely judgement, and judgement is manipulable.

Between roughly 2005 and 2012 several panel banks manipulated
submissions to profit from LIBOR-linked positions and to disguise
their funding costs during the crisis.  Investigations produced
multi-billion-dollar fines and criminal convictions.  The scandal
convinced regulators that a judgement-based benchmark with this
much economic weight was unsustainable.

\begin{historybox}[title={LIBOR scandal 2012}]
Investigations in 2011--2013 revealed coordinated manipulation
of LIBOR submissions at several panel banks.  Emails between
traders and submitters (``an extra basis point would be much
appreciated'') became enforcement exhibits.  Cross-jurisdiction
fines reached roughly \$9~billion; several individuals were
convicted.  In 2017 the UK FCA announced it would no longer
compel LIBOR submissions after end-2021, triggering the global
transition.  Most tenors ceased publication at end-2021; the
final USD settings ceased on 30 June 2023.
\end{historybox}

\subsection{The replacement: risk-free rates}

Each LIBOR-currency regulator designated a \textbf{risk-free
rate (RFR)} as successor.  All RFRs share two features: they are
\emph{overnight}, and based on \emph{observable transactions}
rather than submissions.

\begin{itemize}
  \item \textbf{SOFR} (Secured Overnight Financing Rate):
    USD.  A transaction-based overnight rate backed by US
    Treasury repo volume.  Published by the Federal Reserve
    Bank of New York.
  \item \textbf{SONIA} (Sterling Overnight Index Average):
    GBP.  A transaction-based overnight rate from unsecured
    sterling deposits.  Published by the Bank of England.
  \item \textbf{ESTR} (Euro Short-Term Rate):
    EUR.  Published by the ECB, based on unsecured overnight
    borrowing by euro-area banks.
  \item \textbf{TONA} (Tokyo Overnight Average Rate): JPY.
    Published by the Bank of Japan, unsecured overnight call
    market.
  \item \textbf{SARON} (Swiss Average Rate Overnight):
    CHF, secured, published by SIX Swiss Exchange.
\end{itemize}

\emph{Secured} (SOFR, SARON) versus \emph{unsecured} (SONIA,
ESTR, TONA) is a jurisdictional judgement.  US regulators
concluded that the Treasury repo market was the deepest USD
overnight market; UK, EU, and Japanese regulators preferred
unsecured overnight markets that better mirrored LIBOR.  The
consequence: USD SOFR contains no bank credit component, while
SONIA and ESTR contain a small one.  Basis markets exist between
each RFR and the defunct LIBOR curves.

\subsection{Backward-looking versus forward-looking}

RFRs are \emph{backward-looking}; LIBOR was
\emph{forward-looking}.

LIBOR was a term rate: 3-month USD LIBOR fixed at the start of a
3-month interest period, so the borrower knew the interest amount
on day one.  The rate incorporated a 3-month forward view of
unsecured bank funding.

RFRs are overnight.  A 3-month rate is obtained by compounding
overnight observations daily, so it is known only at the
\emph{end} of the period.  This is operationally different:
treasurers, lenders, and borrowers who planned around knowing
their interest payment in advance had to adapt cashflow systems,
invoicing, and debt covenants to a rate fixed in arrears.  Term
RFR products exist (Term SOFR, Term SONIA), computed from futures
or OIS, but compounded-in-arrears dominates derivatives markets.

\subsection{Fallback language}

Trillions of notional in legacy contracts referenced LIBOR and
matured past cessation.  Bilateral rewrites were infeasible.  The
industry solution, driven by ISDA (International Swaps and
Derivatives Association), was a \textbf{fallback protocol}: a
standard contractual amendment which, once signed, rolls every
in-scope LIBOR reference between two adhering parties to an
RFR-based fallback at cessation.

The fallback rate has two components: a compounded-in-arrears RFR
plus a \textbf{credit-adjustment spread}, set as the historical
median spread between the RFR and the relevant LIBOR tenor over a
five-year lookback window.  The spread is a one-time constant,
fixed at cessation announcement.  Its purpose: ensure the
fallback rate is economically equivalent to LIBOR at transition,
not a step change in coupon.

Loan markets took a slightly different path with contract-specific
``hardwired'' or ``amendment'' fallback language, but the
derivatives market was dominated by the ISDA protocol.

\begin{keyfactbox}[LIBOR to SOFR essentials]
\begin{itemize}
  \item LIBOR: submission-based, forward-looking, unsecured
    term rate.  Ceased end-2021 (most tenors) and
    30 June 2023 (remaining USD settings).
  \item Successor RFRs: SOFR (USD, secured), SONIA (GBP),
    ESTR (EUR), TONA (JPY), SARON (CHF).  All are
    transaction-based, overnight, and backward-looking.
  \item ISDA fallback protocol: compounded RFR plus fixed
    credit-adjustment spread, applied automatically to
    adhering parties at cessation.
  \item Basis markets between RFRs and ex-LIBOR curves
    persist because USD SOFR is secured and the old USD
    LIBOR was unsecured; small but not zero basis everywhere.
\end{itemize}
See \cref{ch:fixed_income_foundations} for the curve-construction
mechanics of SOFR discounting and \cref{ch:rates_derivatives}
for SOFR OIS pricing.
\end{keyfactbox}

\begin{intuitionbox}
LIBOR's problem: it tried to measure something real (bank funding
cost) using a method (quotes) that stopped being tethered to
reality after 2008.  RFRs fix this by measuring overnight markets
that still have volume, at the cost of losing LIBOR's
forward-looking term structure.  Term RFRs and fallback spreads
are the industry's compensation for that loss.
\end{intuitionbox}


%% ================================================================
\section{How regulation touches a Prosperity algorithm}
\label{sec:prosp_reg}
%% ================================================================

\begin{prosperitybox}[title={Rules before code}]
Prosperity has no regulatory layer: no capital requirement, no
stress test, no transaction report, no best-execution
obligation.  A competitor printing sixty thousand orders in a
round is bounded only by the matching engine.

The Prosperity practice that \emph{does} mirror the regulated
world is ``know the rules before you code.''  Each round
publishes a position limit per product, enforced by the engine.
Each round imposes a per-round order-count limit; exceed it and
the algorithm fails to submit.  Tariffs, import penalties, and
conversion constraints (\cref{ch:algo_rounds}) are local
analogues of regulatory constraints.  A strategy optimal in
backtest but violating a rule is worth zero; the first hour of
each round is spent ensuring the algorithm lives inside the
rulebook.

The transferable habit: before any simulation, write the hard
constraints in the same file as the pricing logic.  In
Prosperity those are position and order limits.  On a Goldman
desk they are Pillar~1 capital, FRTB NMRF lists, CCAR stress
sensitivities, MiFID best-execution evidence.  Same activity,
different rulebook.
\end{prosperitybox}


%% ================================================================
\section{How regulation touches a Goldman Strat}
\label{sec:gs_reg}
%% ================================================================

\begin{gsbox}[title={Regulation in a Strat's calendar}]
A Strat in Securities Engineering inherits regulatory calendars
as a background fact of life.

\textbf{FRTB IMA approval is a multi-year project.}  A desk
seeking IMA status submits an application, runs its model in
parallel with SA for at least a year to build a track record,
defends backtesting to the regulator, then continues defending
modellability and NMRF classifications indefinitely.  Strats on
IMA desks spend meaningful time on regulator-facing artefacts:
NMRF evidence files, PLA (profit-and-loss attribution) reports,
backtesting diagnostics.  A desk failing backtesting three times
in a 12-month window is de-approved and falls back to SA until
it re-applies.

\textbf{LIBOR to SOFR was a five-year coordinated push.}  From
roughly 2018, every major bank's rates franchise ran an
enterprise-wide LIBOR transition programme touching trading
systems, pricing libraries, risk engines, client documentation,
and legal.  Discount curves switched from LIBOR to OIS (SOFR,
SONIA, ESTR) discounting; every hard-coded LIBOR pricer was
audited and rewritten; ISDA fallback adherence was coordinated
firm-wide; non-ISDA client contracts were renegotiated
bilaterally.  The transition is not fully finished: legacy
contracts with unusual fallback language remain, basis risk
between term and compounded-in-arrears RFRs is live, and
cross-currency basis swaps carry LIBOR-era residuals.

\textbf{CCAR is an annual drumbeat.}  In Q1 the Fed scenarios
land; February through April the firm's trading-book loss
projections are produced, reconciled, and submitted; June the
Fed publishes results.  During submission season, market-risk
Strats spend significant time mapping macro-scenario variables
into desk-specific shocks, re-running pricers, and explaining
reconciliation differences with DFAST.

\textbf{MiFID II is an ongoing reporting tax.}  Every EU trade
produces a T+1 regulatory report from post-trade infrastructure.
Strats rarely own the reporting pipeline but feel it through
required trade-metadata fields the pricing system must emit:
LEIs, decision-maker codes, algo flags, best-execution evidence.

The regulatory apparatus is slow-moving and relentless.  The
junior Strat who internalises the acronyms in their first year
spends the rest of the career not relearning them.
\end{gsbox}


%% ================================================================
\section{Key facts to memorise}
\label{sec:ch37_key}
%% ================================================================

\begin{keyfactbox}[Chapter summary]
\begin{itemize}
  \item \textbf{Basel III} (2010) is the international
    capital-adequacy framework.  Three pillars: minimum
    capital, supervisory review, market discipline.  Key
    ratios include CET1, Tier~1 $\geq 6\%$, total capital
    $\geq 8\%$, leverage $\geq 3\%$, and the liquidity ratios
    LCR and NSFR.
  \item \textbf{FRTB} (finalised 2019) is the Basel
    market-risk framework.  Replaces VaR at $99\%$ 10-day
    with \textbf{Expected Shortfall at $97.5\%$}, introduces
    per-desk IMA approval versus the standardised approach,
    adds liquidity horizons, and penalises non-modellable
    risk factors (NMRFs) punitively.
  \item \textbf{CCAR} is the Fed's annual stress test for US
    BHCs above roughly \$100 billion.  Macro scenarios drive
    nine-quarter projections; trading-book stressed P\&L
    feeds the Stress Capital Buffer.
  \item \textbf{MiFID II} (EU, 2018) is about market
    structure: pre- and post-trade transparency, best
    execution, research unbundling, T+1 transaction
    reporting, and the Systematic Internaliser regime.
  \item \textbf{LIBOR ceased} at end-2021 for most tenors and
    30 June 2023 for the final USD settings, driven by the
    2012 manipulation scandal.  Replaced by backward-looking
    transaction-based RFRs: SOFR, SONIA, ESTR,
    TONA, SARON.
  \item The \textbf{ISDA fallback protocol} rolled legacy
    LIBOR derivative contracts to compounded RFR plus a
    fixed credit-adjustment spread at cessation.
  \item RFRs are overnight and backward-looking; LIBOR was
    term and forward-looking.  Interest is known only at
    period end under the new convention.
  \item The Strat's daily rhythm is shaped by regulation:
    FRTB modellability, CCAR submission season, MiFID
    reporting metadata, and the long tail of the LIBOR
    transition.
\end{itemize}
\end{keyfactbox}
